\section{case studies}

\subsection{The "Unwind" Musical Biofeedback System}

The "Unwind" system, developed by Yu et al. (2018), employs a special approach to mapping the heart rate to music. It uses nature soundscapes combined with sedative music. Sedative music is used in music therapy to effectively decrease heart rate and blood pressure. It has a low tempo and a repetitive structure\cite{yu2018unwind}.

The system uses the heart rate variability to provide feedback on the user's breathing patterns. The heart rate accelerates during inhalation and decelerates during exhalation. The framework maps these fluctuations directly to a wind sounscape. The wind volume is reduced during inhalation and during exhalation the wind becomes louder. The user should adapt his breahting pattern in a way that a the wind sound is like a smooth, consistent sinve wave. This promotes consistent respiration\cite{yu2018unwind}.

The wind provides immediate feedback on the breathing pattern. The overall relaxation is mapped with a special form of heart rate variability. The system calculates the variability with a sliding window of 16 heartbeats. The organic soundscape is based on the intensity of the Heart Rate Variability. If the variability is low, the environment becomes busy and crowded with bird calls and noises. On the other hand, if the variability increases, the soundscape becomes simpler and more transparent. The frequency of animal sounds decreases and the overall volume drops. This rewards the person with acoustic clarity\cite{yu2018unwind}.

The system was validated with 40 young adults. A significant decrease in heart rate was found in the music assisted conditions, and the reduction in respiration rate was significantly greater with biofeedback than without. Skin conductance responses, which indicate sympathetic nervous system activity, decreased markedly during the relaxation sessions. In follow-up interviews, 70\% of participants reported that the nature sounds allowed them to quickly visualize calming natural scenes, aiding their focus and mood\cite{yu2018unwind}.


\subsection{Technology-Mediated Mindfullness}




